% TEMPLATE-MILC-v3.tex - plantilla maestra UNICA de las guias MILC v3.
%
% La usan las tres familias de guias, cada una con su driver:
%   · Grados 6-11  -> scripts/build-guias-g11.py
%   · Bebras       -> scripts/build-guias-bebras.py
%   · Semillero    -> scripts/build-guias-semillero.py
%
% Antes habia tres copias casi identicas (~400 lineas iguales, 6 distintas):
% cada mejora habia que aplicarla tres veces, y en la practica no se hacia
% (el motor de recursos vivia solo en dos de ellas). Lo que varia entre
% familias esta parametrizado en 6 placeholders que cada driver llena
% (se nombran aqui SIN su sintaxis real: el builder reemplaza por texto plano,
% tambien dentro de los comentarios, y un valor multilinea rompe el preambulo):
%
%   HEADER_IZQ        encabezado izquierdo de pagina
%   HEADER_DER        encabezado derecho de pagina
%   PORTADA_ETIQUETA  rotulo sobre el numero grande (GRADO/MOMENTO/MODULO)
%   PORTADA_SERIE     linea de serie de la portada
%   PORTADA_ID        identificador de la guia en la portada
%   PORTADA_CIRCULO   el circulito de grado (°), o vacio si no aplica
%
% El resto de claves las documenta content/guias/_SCHEMA.md. El script valida
% que no queden placeholders sin reemplazar antes de escribir el archivo final.

\documentclass[11pt,letterpaper]{article}
\usepackage[margin=1.75cm]{geometry}
\usepackage[table]{xcolor}
\usepackage{fontspec}
\usepackage[spanish]{babel}
\usepackage{tikz}
% Librerías TikZ para los diagramas de las guías (bloque `recursos:`):
%  · babel  → babel en español vuelve activos caracteres como «>», y TikZ los usa
%             en sus opciones (>=stealth, ->). Sin esta librería, cualquier
%             diagrama con flechas revienta con «\language@active@arg> has an
%             extra }». Debe ir ANTES de que se dibuje nada.
%  · positioning        → sintaxis «below left=2pt and 3pt».
%  · decorations.pathreplacing → llaves ({brace}) para señalar rangos.
\usetikzlibrary{babel,positioning,decorations.pathreplacing}
\usepackage{graphicx}
% Circuitos: el trabajo de electrónica se hace hoy con micro:bit y sus sensores
% integrados (MakeCode, sin cableado), así que no se necesitan esquemáticos.
% Si algún día entran montajes con protoboard, basta descomentar esta línea:
% los diagramas .tex/.tikz declarados en `recursos:` ya se incrustan solos.
% \usepackage{circuitikz}
\usepackage[most]{tcolorbox}
\usepackage{tabularx}
\usepackage{array}
\usepackage{fancyhdr}
\usepackage{titlesec}
\usepackage{ragged2e}
\usepackage{enumitem}
\usepackage{microtype}

% Helvetica Neue no trae la flecha U+2192, asi que cada «→» escrito literalmente
% en el YAML se compilaba a NADA, en silencio (462 flechas en 61 guias: rutas del
% tipo «paleta → tipografia → lockup» salian rotas). Mapeamos el caracter a su
% version matematica, que si tiene glifo. Asi el autor puede seguir escribiendo
% «→» comodamente en el YAML.
\usepackage{newunicodechar}
\newunicodechar{→}{\ensuremath{\rightarrow}}

\setmainfont{Helvetica Neue}
\setsansfont{Helvetica Neue}
\setlength{\parindent}{0pt}
\setlength{\parskip}{6pt}
\hyphenpenalty=200
\exhyphenpenalty=200
\pretolerance=400
\tolerance=2000
\setlength{\emergencystretch}{1em}
\renewcommand{\arraystretch}{1.25}
\setlist{leftmargin=5mm,itemsep=1mm,topsep=1mm}
\newcolumntype{Y}{>{\RaggedRight\arraybackslash}X}

\definecolor{milcMagenta}{HTML}{E6007E}
\definecolor{milcVino}{HTML}{5A0038}
\definecolor{milcTurquesa}{HTML}{0093A5}
\definecolor{milcVerde}{HTML}{58A923}
\definecolor{milcMostaza}{HTML}{E5B400}
\definecolor{milcOcre}{HTML}{B86600}
\definecolor{milcNegro}{HTML}{241816}
\definecolor{milcGris}{HTML}{F7F7F4}
\definecolor{milcLinea}{HTML}{E8E2DD}
\definecolor{milcGradePrimary}{HTML}{EA580C}
\definecolor{milcGradeDark}{HTML}{9A3412}
\definecolor{milcGradeSoft}{HTML}{FFEDD5}
\definecolor{milcGradeAccent}{HTML}{CFE7FF}
\definecolor{milcGradeText}{HTML}{FFFFFF}
\color{milcNegro}

\pagestyle{fancy}
\fancyhf{}
\lhead{\small Semillero de Investigación · Robótica y sistemas embebidos}
\rhead{\small Módulo 2 · Sistematización · MILC v3}
\cfoot{\small\thepage}
\renewcommand{\headrulewidth}{0pt}

\newtcolorbox{titlebox}[1]{
  enhanced,colback=#1,colframe=#1,arc=7mm,boxrule=0pt,
  left=7mm,right=7mm,top=3mm,bottom=3mm,
  before skip=8pt,after skip=8pt,
  fontupper=\sffamily\bfseries\Large\color{white}
}

\newtcolorbox{softbox}[3]{
  enhanced,colback=#2,colframe=#1,borderline west={4pt}{0pt}{#1},
  arc=2mm,boxrule=.4pt,left=4mm,right=4mm,top=3mm,bottom=3mm,
  before skip=6pt,after skip=8pt,title={#3},coltitle=white,
  colbacktitle=#1,fonttitle=\sffamily\bfseries,fontupper=\RaggedRight,
  attach boxed title to top left={xshift=3mm,yshift=-2mm},
  boxed title style={arc=2mm, boxrule=0pt}
}

\newtcolorbox{drawbox}[2]{
  enhanced,colback=white,colframe=#1,arc=4mm,boxrule=1pt,
  left=5mm,right=5mm,top=4mm,bottom=4mm,before skip=8pt,after skip=8pt,
  title={#2},coltitle=white,colbacktitle=#1,fonttitle=\sffamily\bfseries,
  fontupper=\RaggedRight,
  attach boxed title to top left={xshift=4mm,yshift=-2mm},
  boxed title style={arc=3mm, boxrule=0pt}
}

\newtcolorbox{infoband}[1]{
  enhanced,colback=#1!8,colframe=#1!50,arc=2mm,boxrule=.5pt,
  left=4mm,right=4mm,top=2mm,bottom=2mm,before skip=4pt,after skip=4pt,
  fontupper=\small\RaggedRight
}

% Puente narrativo entre secciones (contrato editorial Conéctate).
% Da continuidad: "Acabas de X. Ahora vas a Y porque Z."
% Visualmente: banda izquierda en color del grado, texto en itálica pequeña.
\newtcolorbox{puentebox}{
  enhanced,colback=milcGradeSoft,colframe=milcGradeDark,
  borderline west={3pt}{0pt}{milcGradeDark},
  arc=0mm,boxrule=0pt,left=5mm,right=4mm,top=2.5mm,bottom=2.5mm,
  before skip=4pt,after skip=8pt,
  fontupper=\itshape\small\color{milcNegro}\RaggedRight
}
% La flecha va en modo matematico a proposito: Helvetica Neue NO tiene el glifo
% U+2192, asi que \textrightarrow se compilaba a NADA («Missing character: There
% is no → in font Helvetica Neue») y el puente salia sin flecha. En modo
% matematico la flecha viene de la fuente matematica y si se dibuja.
\newcommand{\puente}[1]{%
  \begin{puentebox}%
    \textcolor{milcGradeDark}{$\rightarrow$}\hspace{2mm}#1%
  \end{puentebox}%
}

\newcommand{\linead}[1]{\noindent\color{milcLinea}\rule{#1}{0.5pt}\color{milcNegro}}
\newcommand{\respuesta}[1]{\par\vspace{2mm}\foreach \n in {1,...,#1}{\noindent\color{milcLinea}\rule{\linewidth}{0.5pt}\par\vspace{5mm}}\color{milcNegro}}
\newcommand{\checkbox}{\textcolor{milcGradeDark}{\fbox{\phantom{x}}}\hspace{5pt}}

\newcommand{\phasecard}[4]{
\begin{tcolorbox}[enhanced,colback=#3,colframe=#2,arc=3mm,boxrule=.5pt,height=3.05cm,valign=center,halign=center,left=1.5mm,right=1.5mm,top=2mm,bottom=2mm]
{\sffamily\bfseries\fontsize{22}{24}\selectfont\textcolor{#2}{#1}}\par
{\sffamily\bfseries\small #4}
\end{tcolorbox}
}

% ── Recursos visuales de la guía ──────────────────────────────────────
% Declarados en el bloque `recursos:` del YAML y emitidos por el builder.
% \guiaFigura   → imagen rasterizada o PDF (foto, carta celeste, montaje).
% \guiaDiagrama → fuente TikZ/circuitikz (.tex) incrustada como vector:
%                 es la vía para circuitos y esquemáticos editables.
% #1 = ruta relativa al .tex · #2 = pie (puede ir vacío)
\newcommand{\guiaFigura}[2]{%
\begin{center}
\includegraphics[width=0.88\linewidth,height=0.38\textheight,keepaspectratio]{#1}\\[1.5mm]
{\footnotesize\itshape\textcolor{milcNegro!75}{#2}}
\end{center}
}

\newcommand{\guiaDiagrama}[2]{%
\begin{center}
\input{#1}\\[1.5mm]
{\footnotesize\itshape\textcolor{milcNegro!75}{#2}}
\end{center}
}

\begin{document}
\thispagestyle{empty}
\begin{tikzpicture}[remember picture,overlay]
  \fill[white] (current page.south west) rectangle (current page.north east);
  \path[fill=milcGradePrimary,rounded corners=16mm]
    ([xshift=.55cm,yshift=-.55cm]current page.north west) rectangle
    ([xshift=-.55cm,yshift=3.65cm]current page.south east);

  \node[anchor=north west,text=milcGradeText] at ([xshift=2.0cm,yshift=-1.85cm]current page.north west)
    {\sffamily\bfseries\fontsize{14}{16}\selectfont MÓDULO};
  \node[anchor=north west,text=milcGradeText,opacity=.82] at ([xshift=2.0cm,yshift=-2.75cm]current page.north west)
    {\sffamily\bfseries\fontsize{9}{11}\selectfont SEMILLERO DE INVESTIGACIÓN · CONECTATE};

  \begin{scope}[shift={([xshift=-3.05cm,yshift=-2.55cm]current page.north east)}]
    \fill[white,opacity=.80] (-.62,-.52) rectangle (.62,.52);
    \draw[milcGradeText,line width=1pt] (-.62,-.52) rectangle (.62,.52);
    \fill[milcGradeDark] (-.42,-.38) rectangle (-.22,.06);
    \fill[milcGradeAccent] (-.10,-.38) rectangle (.10,.24);
    \fill[milcGradePrimary!60!black] (.22,-.38) rectangle (.42,.38);
  \end{scope}

  \node[anchor=west,text=milcGradeText] at ([xshift=1.9cm,yshift=-12.85cm]current page.north west)
    {\sffamily\bfseries\fontsize{124}{124}\selectfont 2};
  % El circulito de grado (°) solo aplica a las guias de grado («11°»); en
  % Bebras y Semillero el numero grande es un momento/modulo, no un grado.
  

  \node[anchor=west,text=milcGradeText,text width=8.7cm,align=left] at ([xshift=10.35cm,yshift=-11.6cm]current page.north west)
    {
     {\sffamily\bfseries\fontsize{19}{22}\selectfont De la idea al pseudocódigo ---\\de la cera perdida\\de los Quimbaya\\al diagrama de flujo}\\[4mm]
     {\sffamily\bfseries\fontsize{23}{26}\selectfont Robótica y sistemas embebidos}\\[2mm]
     {\sffamily\fontsize{13}{16}\selectfont Módulo 2 · Fase MILC: Sistematización · 3 h de clase}\\[5mm]
     {\sffamily\bfseries\fontsize{11}{13}\selectfont Institución Educativa Sor María Juliana}\\[1mm]
     {\sffamily\fontsize{10}{12}\selectfont Autor: Dr. Álvaro Cárdenas Orozco}
    };

  \path[fill=milcGradeSoft,rounded corners=7mm]
    ([xshift=1.35cm,yshift=.85cm]current page.south west) rectangle
    ([xshift=-1.35cm,yshift=4.25cm]current page.south east);
  \draw[milcGradeDark,line width=.65pt,rounded corners=7mm]
    ([xshift=1.35cm,yshift=.85cm]current page.south west) rectangle
    ([xshift=-1.35cm,yshift=4.25cm]current page.south east);
  \node[anchor=center,text=milcNegro,text width=17.0cm,align=center] at ([yshift=2.55cm]current page.south)
    {\sffamily\fontsize{10}{12}\selectfont Metodología MILC · Apertura ancestral · Pensamiento computacional · Triángulo Dussel-Estoicismo-Floridi\\
    \textcolor{milcGradeDark}{\bfseries Producto:} La \textbf{lógica de tu aparato diseñada antes de programar}: el \textbf{pseudocódigo} (la lógica en palabras: percibir → decidir con una regla → actuar) y el \textbf{diagrama de flujo} (la misma lógica dibujada con sus figuras). Más tu \textbf{tabla de umbrales} ---qué valor separa cada decisión y \textbf{cómo lo vas a medir} para justificarlo--- y la \textbf{lista de materiales} de tu prototipo. Todo listo para que en el módulo 3 los bloques de MakeCode solo lo \emph{traduzcan}.};
\end{tikzpicture}
\mbox{}
\newpage

\section*{Apertura: De la idea al pseudocódigo --- de la cera perdida de los Quimbaya al diagrama de flujo}

\begin{softbox}{milcGradeDark}{milcGradeSoft}{Saber ancestral}
\textbf{En las vertientes del Norte del Valle, donde hoy está Cartago, los Quimbaya fueron maestros orfebres.} Sus piezas más finas ---el poporo, las figuras huecas--- nacían de una técnica exigente: la \textbf{fundición a la cera perdida} con núcleo (Banco de la República). Y esa técnica es, sin saberlo, un \textbf{algoritmo perfecto}. El maestro modelaba la figura en \textbf{cera de abejas} sobre un núcleo de arcilla y carbón; la cubría con arcilla dejando un canal; \textbf{calentaba} el molde hasta que la cera se derretía y salía, dejando su forma hueca; entonces \textbf{vertía} la tumbaga (oro y cobre) fundida en ese vacío; y al enfriar, \textbf{rompía} el molde para revelar la pieza. Cada paso dependía del anterior, y el orden era \textbf{absoluto}: si vertía el oro sin haber sacado la cera, la pieza se perdía. Por eso el orfebre \textbf{planeaba toda la secuencia con precisión antes de empezar}: porque una vez vertido el oro, ya no había corrección posible. Programar un aparato es igual: primero se diseña la lógica exacta ---qué percibe, qué decide, qué hace, y en qué orden---, y solo después se ejecuta. \textbf{El Quimbaya no improvisaba con el oro; tú no improvisas con los bloques.}
\end{softbox}

\begin{softbox}{milcTurquesa}{milcGris}{Saber contemporáneo}
Entre la idea y el programa hay un puente que los principiantes se saltan ---y por eso se pierden: \textbf{el diseño de la lógica}. Sistematizar es escribir \emph{cómo} va a pensar tu aparato, con dos herramientas. \textbf{(1) El pseudocódigo:} la lógica en \textbf{palabras ordenadas}, sin bloques todavía. \emph{``Por siempre: leer el sensor. Si el valor es menor que el umbral, hacer A; si no, hacer B''}. Es el plano en lenguaje humano. \textbf{(2) El diagrama de flujo:} la misma lógica \textbf{dibujada}, con cuatro figuras ---óvalo (inicio/fin), rectángulo (una acción), rombo (una pregunta sí/no) y flechas (el camino). Ver la lógica dibujada delata los huecos que las palabras esconden. En el corazón de casi todo está el \textbf{umbral}: el sensor entrega un \emph{número}, y el umbral es el valor de corte que lo vuelve \textbf{decisión} (\emph{¿la luz es menor que 30?}). Pero un umbral no se inventa: se \textbf{calibra} midiendo condiciones conocidas ---como hiciste con el fotómetro en Astronomía, como la curandera fijaba «fiebre» con su mano. \textbf{La regla de oro:} corregir la lógica en palabras o en un dibujo cuesta un borrón; corregirla con los bloques ya armados cuesta horas ---y con el oro ya vertido, cuesta la pieza.
\end{softbox}

\begin{softbox}{milcMostaza}{milcGris}{Pregunta-puente}
\textit{El orfebre Quimbaya no vertía el oro hasta tener toda la secuencia clara en la cabeza, porque el error ya no se podía deshacer; \textbf{¿por qué diseñar la lógica de tu aparato ---en pseudocódigo y diagrama--- antes de tocar un bloque te ahorra el «oro perdido»} de programar a ciegas?}
\end{softbox}

\begin{softbox}{milcVerde}{milcGris}{Saber y saber hacer}
Al terminar la sesión: (1) podrás \textbf{identificar} las partes de una lógica (sensor, umbral, decisión, acción) dentro de tu necesidad del módulo 1; (2) sabrás \textbf{explicar} qué es el pseudocódigo, las cuatro figuras del diagrama de flujo y qué es un umbral calibrado; (3) podrás \textbf{analizar} tu propia lógica buscando huecos e incoherencias entre lo que percibe y lo que hace; (4) habrás \textbf{diseñado} el pseudocódigo, el diagrama de flujo y la tabla de umbrales de tu aparato.
\end{softbox}



\small
\begin{tabularx}{\textwidth}{>{\bfseries}p{3.0cm}Y}
\rowcolor{milcGradePrimary!12}Autor & Dr. Álvaro Cárdenas Orozco\\
DBA & Formula preguntas investigables, produce evidencia con método y comunica hallazgos reconociendo sus límites (investigación estudiantil STEM, transversal 6°--11°)\\
Referentes & micro:bit · MakeCode · Continuidad con la unidad de 8° P2 (guías 8-2-1 a 8-2-10)\\
Producto final & La \textbf{lógica de tu aparato diseñada antes de programar}: el \textbf{pseudocódigo} (la lógica en palabras: percibir → decidir con una regla → actuar) y el \textbf{diagrama de flujo} (la misma lógica dibujada con sus figuras). Más tu \textbf{tabla de umbrales} ---qué valor separa cada decisión y \textbf{cómo lo vas a medir} para justificarlo--- y la \textbf{lista de materiales} de tu prototipo. Todo listo para que en el módulo 3 los bloques de MakeCode solo lo \emph{traduzcan}.\\
\end{tabularx}
\normalsize

\puente{En el módulo 1 escuchaste tu entorno y afinaste una necesidad (percibir → regla → hacer). Hoy le das \textbf{forma exacta}: la conviertes en una lógica que otro podría leer y construir, como el plano del orfebre antes de fundir.}

\begin{titlebox}{milcGradeDark}
Ruta MILC: así vamos a aprender
\end{titlebox}
\begin{tabularx}{\textwidth}{>{\centering\arraybackslash}X>{\centering\arraybackslash}X>{\centering\arraybackslash}X>{\centering\arraybackslash}X}
\phasecard{1}{milcVerde}{milcGris}{Escucha\\[-1mm]\small Observo, escucho y pregunto.} &
\phasecard{2}{milcTurquesa}{milcGris}{Sistematización\\[-1mm]\small Pensamiento computacional.} &
\phasecard{3}{milcMagenta}{milcGris}{Praxis\\[-1mm]\small Construyo y aplico.} &
\phasecard{4}{milcVino}{milcGris}{Evaluación\\[-1mm]\small Reflexión liberadora.}
\end{tabularx}


\puente{Antes de diseñar, vuelve a tu necesidad afinada del módulo 1. Ese es el barro: hoy lo modelas en una lógica precisa.}

\begin{titlebox}{milcVerde}
Fase 1 · Escucha
\end{titlebox}

\begin{softbox}{milcVerde}{milcGris}{Escena para escuchar}
\textbf{Actividad 1 · IDENTIFICA --- Las partes de una lógica} (35 min · individual + grupo). \textbf{Momento A (10 min) --- El caso de la cera perdida.} El profe cuenta la secuencia del orfebre Quimbaya y pregunta: \emph{¿qué pasaba si se saltaba o cambiaba un paso? ¿por qué no podía improvisar?}. \textbf{Momento B (15 min) --- Desarmar una lógica cotidiana.} En grupo, tomen un aparato conocido (una nevera, una puerta automática, una alarma) y sepárenlo en sus piezas lógicas: \emph{¿qué percibe (sensor)? ¿qué número o señal mira? ¿cuál es el umbral (el valor de corte)? ¿qué decide? ¿qué hace (actuador)?}. \textbf{Momento C (10 min) --- Mi necesidad, en piezas.} Cada quien toma su necesidad afinada del módulo 1 y la parte en esas mismas piezas: sensor, umbral, decisión, acción. \textbf{En el cuaderno:} el aparato cotidiano desarmado en piezas lógicas + tu necesidad partida en sensor / umbral / decisión / acción.
\end{softbox}

\checkbox Puedo explicar por qué el orfebre no podía cambiar el orden de los pasos.\\
\checkbox Desarmé un aparato cotidiano en sensor, umbral, decisión y actuador.\\
\checkbox Partí mi propia necesidad en esas cuatro piezas.

\begin{infoband}{milcVerde}


\textbf{Lo que va al cuaderno (Actividad 1 · IDENTIFICA).} Título: «Actividad 1 --- Las partes de una lógica». \textbf{Formato:} un aparato cotidiano desarmado en sus piezas + mi necesidad del módulo 1 partida en sensor / umbral / decisión / acción. \textbf{Extensión:} 8 líneas. \textbf{Sabes que terminaste cuando} puedes nombrar, en tu necesidad, \textbf{cuál es el umbral} (el número que separa una decisión de otra).
\end{infoband}

\puente{Ya tienes tu necesidad. Ahora aprende a darle forma: pseudocódigo, diagrama de flujo y umbrales calibrados ---las herramientas de pensar antes de programar.}

\begin{titlebox}{milcTurquesa}
Fase 2 · Sistematización con pensamiento computacional
\end{titlebox}

\begin{softbox}{milcTurquesa}{milcGris}{Cuatro pilares aplicados al tema}
\textbf{Actividad 2 · EXPLICA --- Pensar antes de programar} (70 min · grupo + individual). Ya tienes tu necesidad en piezas. Ahora aprende a \textbf{ensamblarlas en una lógica}: el pseudocódigo, el diagrama de flujo con sus cuatro figuras, y cómo se \textbf{calibra} un umbral en vez de inventarlo.
\end{softbox}

\small
\begin{tabularx}{\textwidth}{>{\bfseries\RaggedRight\arraybackslash}p{3.6cm}Y}
\rowcolor{milcTurquesa!90}\color{white}Pilar & \color{white}\bfseries Aplicación al tema\\
1. Descomposición & \textbf{Pilar 1 --- El pseudocódigo: la lógica en palabras.} Antes de un solo bloque, se escribe \emph{cómo pensará} el aparato, en frases ordenadas que cualquiera entiende. La columna vertebral casi siempre es la misma: \textbf{un bucle} (\emph{``por siempre''}, porque un sistema vigila sin parar) que adentro \textbf{lee} el sensor y \textbf{decide} con una regla. Ejemplo: \emph{``Por siempre: leer la luz. Si la luz es menor que 30, mostrar la luna. Si no, mostrar el sol''}. \textbf{Léelo en voz alta:} si suena coherente de principio a fin, vas bien; si te trabas, ahí hay un hueco. \emph{El pseudocódigo es barato de corregir ---es solo texto--- y por eso es donde deben morir los errores.}\\
2. Reconocimiento de patrones & \textbf{Pilar 2 --- El diagrama de flujo: la lógica dibujada, cuatro figuras.} Lo que las palabras esconden, el dibujo lo delata. Un diagrama de flujo usa \textbf{cuatro figuras}: el \textbf{óvalo} marca dónde \emph{empieza y termina}; el \textbf{rectángulo} es \emph{una acción} (leer, mostrar); el \textbf{rombo} es \emph{una pregunta} de sí/no (\emph{¿luz $<$ 30?}), y de él salen \textbf{dos} caminos; las \textbf{flechas} muestran por dónde va el proceso. \textbf{La prueba de fuego:} sigue las flechas con el dedo; si llegas a un rectángulo del que no sale flecha, o a un rombo con una sola salida, \textbf{tienes un hueco}. Dibujar obliga a decidir lo que el texto deja vago. Es exactamente lo que aprendiste en 7°, ahora al servicio de tu aparato.\\
3. Abstracción & \textbf{Pilar 3 --- El umbral: el número que se vuelve decisión.} El sensor no dice «oscuro»: dice \emph{42}. El \textbf{umbral} es el valor de corte que convierte ese número en una decisión (\emph{si es menor que 30, es «oscuro»}). Sin umbral no hay decisión; con un umbral inventado, la decisión es un capricho. Por eso el umbral se \textbf{calibra}: se \textbf{mide} el sensor en condiciones \emph{conocidas} (tapado = oscuro, bajo la luz = claro) y se elige el corte con esos números en la mano. \emph{Ya lo hiciste con el fotómetro en Astronomía, y es lo que hacía la curandera del Pacífico al fijar «fiebre leve» con su propia mano como patrón.} Un umbral calibrado se puede \textbf{justificar}; uno inventado, solo defender.\\
4. Algoritmo conceptual & \textbf{Pilar 4 --- Coherencia y materiales: que todo encaje.} Una lógica sirve cuando lo que \textbf{percibe} y lo que \textbf{hace} encajan sin huecos: cada pregunta tiene sus dos respuestas, cada camino termina en una acción, ninguna decisión queda colgando. Revísalo como el orfebre revisaba su secuencia: \emph{de principio a fin, sin saltos}. Y ---detalle práctico--- haz tu \textbf{lista de materiales}: para la mayoría de necesidades basta \textbf{un micro:bit} (sus sensores vienen adentro); algunas piden algo más (una pila, un cable, un sensor externo). Saber qué necesitas \emph{antes} de construir es parte de planear como el que funde: no se enciende el horno sin tener todo listo.\\
\end{tabularx}
\normalsize

\begin{drawbox}{milcTurquesa}{Pseudocódigo + diagrama de flujo + umbral + coherencia}
\textbf{Pseudocódigo:} la lógica en palabras. Bucle «por siempre» → leer sensor → si (regla) hacer A, si no hacer B. Barato de corregir. \\[1mm]
\textbf{Diagrama de flujo:} óvalo (inicio/fin) · rectángulo (acción) · rombo (pregunta sí/no, dos salidas) · flechas (camino). \\[1mm]
\textbf{Umbral:} el valor de corte que vuelve el número una decisión. Se calibra midiendo condiciones conocidas, no se inventa. \\[1mm]
\textbf{Coherencia:} cada pregunta con sus dos salidas, cada camino termina en acción. Más la lista de materiales.
\end{drawbox}

\begin{softbox}{milcMostaza}{milcGris}{Errores comunes y su corrección}
\textbf{Errores al diseñar la lógica (y cómo evitarlos).} (1) \emph{Saltar directo a los bloques} --- programar sin plano es fundir sin secuencia; \textbf{escribe el pseudocódigo primero}. (2) \emph{Umbral inventado} --- «pongo 50 porque sí» no se puede justificar; \textbf{calíbralo con medidas}. (3) \emph{Rombo con una sola salida} --- toda pregunta tiene sí \emph{y} no; si falta un camino, hay un hueco. (4) \emph{Acción sin salida} --- un paso del que no sale flecha deja el proceso colgado. (5) \emph{Percibir sin actuar} --- si la lógica lee pero no hace nada visible, no es un sistema. (6) \emph{Diagrama que solo tú entiendes} --- si un compañero no puede seguir las flechas sin tu ayuda, todavía no está claro; el plano es para que \emph{otro} lo lea.
\end{softbox}

\begin{infoband}{milcTurquesa}
\guiaDiagrama{assets/robotica-2/diagrama-de-flujo.tex}{El diagrama de flujo dibuja la lógica con cuatro figuras: óvalo (inicio/fin), rectángulo (acción), rombo (pregunta sí/no) y flechas (el camino). Si se lee claro aquí, los bloques solo lo traducen.}

\textbf{Lo que va al cuaderno (Actividad 2 · EXPLICA).} Título: «Actividad 2 --- Pensar antes de programar». \textbf{Formato:} 3 bloques. Bloque A: un pseudocódigo de ejemplo (bucle + lectura + decisión). Bloque B: las cuatro figuras del diagrama de flujo, cada una con lo que significa. Bloque C: la anatomía de un umbral calibrado (mido tapado = \_\_, mido con luz = \_\_, pongo el corte en \_\_). \textbf{Extensión:} 1 hoja. \textbf{Sabes que terminaste cuando} podrías explicar \textbf{por qué un umbral se mide y no se inventa}.
\end{infoband}

\puente{Ya conoces las herramientas. Es hora de usarlas con \emph{tu} necesidad: escribir su lógica, dibujarla y buscarle los huecos antes de que cuesten caro.}

\begin{titlebox}{milcMagenta}
Fase 3 · Praxis con pensamiento computacional
\end{titlebox}

\begin{softbox}{milcMagenta}{milcGris}{Construyendo el producto con los cuatro pilares}
\textbf{Actividad 3 · ANALIZA --- Diseña y depura la lógica de tu aparato} (65 min · individual + parejas). Hora de hacer el plano. Vas a escribir el \textbf{pseudocódigo} de tu necesidad, dibujarlo como \textbf{diagrama de flujo}, fijar tus \textbf{umbrales} y ---clave--- \textbf{buscarle los huecos} antes de que cuesten caro en el módulo 3.
\end{softbox}

\small
\begin{tabularx}{\textwidth}{>{\bfseries\RaggedRight\arraybackslash}p{3.6cm}Y}
\rowcolor{milcMagenta!90}\color{white}Pilar & \color{white}\bfseries Aplicación al producto\\
Descomposición & \textbf{Paso 1 --- Escribe el pseudocódigo (15 min).} Toma tu necesidad afinada (percibir → regla → hacer) y escríbela como lógica en palabras, empezando por el bucle: \emph{``Por siempre: leer \_\_\_\_ (sensor). Si \_\_\_\_ (regla con umbral), hacer \_\_\_\_ (acción). Si no, hacer \_\_\_\_''}. Usa tantas preguntas como necesite tu idea (una alerta de 3 estados necesita dos umbrales). \textbf{Léelo en voz alta} de principio a fin: donde te trabes, corrige ---aquí es barato.\\
Patrones & \textbf{Paso 2 --- Dibuja el diagrama de flujo (20 min).} Traduce tu pseudocódigo a las cuatro figuras: \textbf{óvalo} para empezar, \textbf{rectángulos} para las acciones (leer, mostrar), \textbf{rombos} para las preguntas del umbral (con sus \emph{dos} flechas: sí / no), y \textbf{flechas} para el camino, cerrando el bucle al inicio. Hazlo grande y limpio. Este dibujo es tu \textbf{cera}: lo que aquí quede bien, luego se funde sin sorpresas.\\
Abstracción & \textbf{Paso 3 --- Fija y justifica tus umbrales (10 min).} Para cada rombo, decide el \textbf{valor de corte} y escribe \textbf{cómo lo vas a calibrar}: \emph{``mediré el sensor tapado y con luz, y pondré el umbral en el punto medio''}. Si aún no tienes el micro:bit en mano, deja el número \emph{provisional} y anota el método ---lo confirmarás midiendo en el módulo 3. Un umbral con método vale; uno inventado, no.\\
Algoritmo de armado & \textbf{Paso 4 --- La prueba del compañero (10 min).} Intercambia tu diagrama con tu pareja. Él debe \textbf{seguir las flechas con el dedo} y contarte, sin tu ayuda, qué hace tu aparato en cada caso. \textbf{Cada vez que se pierda, encontraste un hueco:} un rombo sin dos salidas, una acción sin flecha, un camino que no vuelve. Márcalos y arréglalos. \emph{Este es el momento más valioso del módulo: cada hueco tapado hoy es una hora ahorrada mañana.}\\
\end{tabularx}
\normalsize

\begin{drawbox}{milcMagenta}{Antes de dar por bueno tu plano}
\checkbox Escribí el pseudocódigo con su bucle, su lectura y sus decisiones. \\[1mm]
\checkbox Dibujé el diagrama de flujo con las cuatro figuras y el bucle cerrado. \\[1mm]
\checkbox Cada rombo tiene sus dos salidas (sí / no) y cada camino termina en acción. \\[1mm]
\checkbox Cada umbral tiene un valor y un método de calibración escrito. \\[1mm]
\checkbox Mi pareja siguió las flechas sin mi ayuda (o marqué y arreglé los huecos). \\[1mm]
\checkbox Tengo la lista de materiales de mi prototipo.
\end{drawbox}

\begin{softbox}{milcMagenta}{milcGris}{Plantilla de trabajo}
\textbf{Plantilla del plano.} \textit{Pseudocódigo:} por siempre → leer \_\_\_\_ → si \_\_\_\_ (umbral) hacer \_\_\_\_ / si no hacer \_\_\_\_. \\[1mm]
\textit{Diagrama de flujo:} óvalo inicio → rectángulos (acciones) → rombos (preguntas, 2 salidas) → flechas → vuelve al inicio. \\[1mm]
\textit{Umbrales:} umbral 1 = \_\_ (calibro midiendo \_\_\_\_) · umbral 2 = \_\_ (calibro midiendo \_\_\_\_). \\[1mm]
\textit{Materiales:} micro:bit + \_\_\_\_.
\end{softbox}

\begin{infoband}{milcMagenta}


\textbf{Lo que va al cuaderno (Actividad 3 · ANALIZA).} Título: «Actividad 3 --- El plano de mi aparato». \textbf{Formato:} el pseudocódigo + el diagrama de flujo depurado (con los huecos ya marcados y arreglados) + la tabla de umbrales con su método + la lista de materiales. \textbf{Extensión:} 1-2 hojas. \textbf{Sabes que terminaste cuando} otra persona pudo seguir tu diagrama de flujo sin que se lo explicaras.
\end{infoband}

\puente{El producto no es un programa: es su \textbf{plano} ---pseudocódigo, diagrama y umbrales. Con eso, el módulo 3 será traducir, no adivinar.}

\begin{titlebox}{milcMostaza}
Producto de la sesión
\end{titlebox}
\textbf{Plano lógico del aparato — pseudocódigo + diagrama de flujo + umbrales + materiales}.
\begin{softbox}{milcMostaza}{milcGris}{Criterios de calidad}
(1) El pseudocódigo tiene bucle, lectura del sensor y decisiones con umbral. (2) El diagrama de flujo usa las cuatro figuras y cada rombo tiene dos salidas. (3) Ningún camino queda colgado y el bucle se cierra (no hay huecos en la lógica). (4) Cada umbral tiene un valor y un método de calibración justificado (no inventado). (5) Están la lista de materiales y la coherencia entre lo que percibe y lo que hace.
\end{softbox}

\puente{Antes de cerrar, prueba tu diagrama con un compañero: si él no puede seguir la lógica sin que se la expliques, todavía tiene huecos.}

\begin{titlebox}{milcVino}
Fase 4 · Evaluación liberadora — cinco dimensiones
\end{titlebox}

\begin{softbox}{milcVino}{milcGris}{Sentido de la evaluación}
Evaluar no es solo revisar una nota. Es reconocer qué aprendí, qué cambió en mí, qué emociones manejé, cómo me relaciono con la ciudad y la comunidad, y qué le dejaré a quien viene detrás.
\end{softbox}

\textbf{Mi reflexión liberadora en cinco dimensiones.} Completa cada frase iniciada:

\small
\begin{tabularx}{\textwidth}{>{\bfseries\RaggedRight\arraybackslash}p{3.4cm}Y>{\RaggedRight\arraybackslash}p{4.6cm}}
\rowcolor{milcVino}\color{white}Dimensión & \color{white}\bfseries Mi respuesta (frame starter) & \color{white}\bfseries Algo que descubrí\\
1. Personal & Lo que más cambió en mí fue \linead{6.5cm} & \linead{4.2cm}\\
2. Emocional & La emoción más fuerte fue \linead{2.5cm} y la regulé \linead{3cm} & \linead{4.2cm}\\
3. Ciudadana & En democracia esto importa porque \linead{6cm} & \linead{4.2cm}\\
4. Local (Cartago) & En mi barrio sería útil para \linead{6.5cm} & \linead{4.2cm}\\
5. Intergeneracional & A mi abuela/abuelo le diría \linead{6.5cm} & \linead{4.2cm}\\
\end{tabularx}
\normalsize

\begin{infoband}{milcVino}
\textbf{Para la próxima sustentación.} Vuelve a esta tabla en seis meses. Verás que las respuestas cambian — no porque lo aprendido cambie, sino porque tú cambias. La evaluación liberadora es una conversación contigo a través del tiempo.
\end{infoband}


\puente{Cerramos con 3 voces que te ayudan a entender por qué pensar el orden antes de actuar ---y no al revés--- es lo que distingue al oficio de la improvisación.}

\begin{titlebox}{milcGradeDark}
Triángulo de pensamiento · cierre filosófico
\end{titlebox}

\begin{softbox}{milcVino}{milcGris}{Dussel · lente del nosotros}
\textit{``Analéctico quiere indicar el hecho real humano por el que todo hombre, todo grupo o pueblo se sitúa siempre más allá (aná-) del horizonte de la totalidad.''}

Dussel dice que hay saberes que vienen de más allá del sistema establecido. Tu necesidad ---la de tu casa, tu vereda--- suele ser una de esas: no está en ningún catálogo de proyectos. \textbf{Diseñar su lógica es darle forma rigurosa para que exista}: el pseudocódigo y el diagrama convierten una intuición situada en un plano que otro puede leer y construir. \emph{Sistematizar no le quita lo propio a tu idea; le da los huesos para ponerse de pie.}

\textbf{¿Le estoy dando a mi idea una forma clara para que otro la entienda, o la dejo solo en mi cabeza?}
\end{softbox}
\linead{\linewidth}\\[1mm]
\linead{\linewidth}

\begin{softbox}{milcMostaza}{milcGris}{Estoicismo · Epicteto · Enquiridión, 1 (c. 125 d.C.; trad. desde E. Carter) · cuidado interior}
\textit{``Algunas cosas están en nuestro poder y otras no. Están en nuestro poder la opinión, el impulso, el deseo, el rechazo; en una palabra, todo aquello que es obra nuestra.''}

Epicteto separa lo que depende de mí de lo que no. No controlas si el sensor será perfecto ni si el mundo cooperará; \textbf{sí controlas la claridad de tu plan}. Un diagrama sin huecos, un umbral con método, una lógica que otro puede seguir: eso es \emph{obra tuya}, y es donde se gana o se pierde un proyecto. El orfebre no controlaba el fuego, pero sí la secuencia. \textbf{Domina lo que depende de ti ---el diseño--- y el resto se vuelve manejable.}

\textbf{¿Estoy puliendo lo que sí depende de mí (la claridad de mi lógica), o culpando a lo que no?}
\end{softbox}
\linead{\linewidth}\\[1mm]
\linead{\linewidth}

\begin{softbox}{milcTurquesa}{milcGris}{Floridi · lente de la infoesfera}
\textit{``Las tecnologías digitales están re-ontologizando la naturaleza misma de la infoesfera.''}

\emph{Re-ontologizar} es transformar la naturaleza de algo. Tu idea, mientras vive solo en tu cabeza, es vaga: cambia cada vez que la piensas. Escribirla como pseudocódigo y diagrama de flujo \textbf{la vuelve otra cosa}: algo preciso, fijo, ejecutable y compartible. El plano no \emph{describe} tu idea: la \textbf{re-ontologiza} en algo que una máquina puede correr y una persona puede revisar. \textbf{Diseñar la lógica es el acto que convierte un deseo en un artefacto posible.}

\textbf{Mi lógica, ¿ya es precisa y ejecutable, o sigue siendo una idea que cambia cada vez que la cuento?}
\end{softbox}
\linead{\linewidth}\\[1mm]
\linead{\linewidth}

\begin{titlebox}{milcVino}
Compromiso final
\end{titlebox}

\small
\begin{tabularx}{\textwidth}{>{\RaggedRight\arraybackslash}p{3.0cm}YYY}
\rowcolor{milcVino}\color{white}\bfseries Criterio & \color{white}\bfseries Lo logré & \color{white}\bfseries En proceso & \color{white}\bfseries Necesito apoyo\\
Comprensión del tema & Explico ideas centrales con mis palabras. & Reconozco ideas pero falta precisión. & Necesito repasar conceptos base.\\
Pensamiento computacional & Apliqué los 4 pilares al diseñar mi producto. & Apliqué algunos pilares. & Necesito releer la fase 2.\\
Aplicación y producto & Mi producto cumple los criterios y se puede revisar. & Producto funcional con ajustes pendientes. & Aún en construcción.\\
Reflexión 5 dimensiones & Respondí cada dimensión con honestidad. & Respondí algunas con detalle. & Mis respuestas son muy generales.\\
Triángulo de pensamiento & Conecté las 3 voces con mi trabajo real. & Conecté al menos una. & No relaciono las voces con mi proceso.\\
\end{tabularx}
\normalsize

\textbf{Hoy entendí que:}\\[1mm]
\linead{\linewidth}\\[1mm]
\linead{\linewidth}

\textbf{Mi compromiso para la próxima sesión es:}\\[1mm]
\linead{\linewidth}\\[1mm]
\linead{\linewidth}

\begin{softbox}{milcMostaza}{milcGris}{Cita de despedida · Marco Aurelio}
\textit{``El obstáculo en el camino se vuelve el camino. Lo que estorba al camino se vuelve camino.''}\\[1mm]
Cada error de hoy, bien observado, se convierte en pieza del camino que pisarás mañana.
\end{softbox}

\small
\begin{tabularx}{\textwidth}{>{\bfseries}p{3.6cm}Y}
\rowcolor{milcGradeDark!12}Nombre completo: & \linead{10cm}\\
Fecha: & \linead{4cm} \quad Curso/grupo: \linead{4cm}\\
Mi firma: & \linead{9cm}\\
\end{tabularx}
\normalsize

\begin{infoband}{milcGradeDark}
\textbf{Para la próxima sesión.} Dos cosas. \textbf{(1) Pule tu plano:} deja el pseudocódigo, el diagrama de flujo (sin huecos) y la tabla de umbrales listos en una hoja limpia. Haz que \textbf{otra persona} ---de tu casa o tu curso--- siga tu diagrama con el dedo y te cuente qué hace el aparato; si se pierde, arréglalo. \textbf{(2) Consigue tu material:} averigua si tienes acceso a un \textbf{micro:bit} (propio, del colegio, o compartido en grupo) y qué más de tu lista de materiales necesitas. Si no hay aparato físico, no importa: en el módulo 3 usaremos también el \textbf{simulador de MakeCode}, que trae los sensores en pantalla. \textbf{Trae:} tu plano completo y depurado. \emph{En el módulo 3 pasas de la cera al oro: traduces tu diagrama a bloques de MakeCode, mides tus umbrales de verdad y construyes el prototipo. Hoy hiciste lo que el orfebre hacía antes de encender el horno ---y por eso tu fundición no se va a perder.}
\end{infoband}

\end{document}
